\exercisesection{1}

\begin{enumerate}
\def\labelenumi{\arabic{enumi}.}
\tightlist
\item
  Let K be a field; for every subring A of K let L(A) denote the set of local rings \(A_{p}\) of A where p runs through the prime ideals of A (these local rings being identified with subrings of K).
\end{enumerate}

\begin{enumerate}
\def\labelenumi{(\alph{enumi})}
\item
  For a local subring M of K to dominate a ring \(A_{p} \in L(A)\) , it is necessary and sufficient that \(A \subset M\) ; the local ring \(A_{p}\) dominated by M is then unique and corresponds to the prime ideal \(p = m(M) \cap A\) .
\item
  Let M, N be two local subrings of K and P the subring of K generated by \(M \cup N\) . Show that the following conditions are equivalent:
\end{enumerate}

\((\alpha)\) There exists a prime ideal \(\mathfrak{p}\) of \(\mathbf{P}\) such that \(\mathfrak{m}(\mathbf{M}) = \mathfrak{p} \cap \mathbf{M}\) , \(\mathfrak{m}(\mathbf{N}) = \mathfrak{p} \cap \mathbf{N}\) .

(β) The ideal \(\mathfrak{a}\) generated in \(\mathbf{P}\) by \(\mathfrak{m}(\mathbf{M}) \cup \mathfrak{m}(\mathbf{N})\) is distinct from \(\mathbf{P}\) .

(γ) There exists a local subring Q of K dominating both M and N.

If these conditions are satisfied, M and N are called related.

\begin{enumerate}
\def\labelenumi{(\alph{enumi})}
\setcounter{enumi}{2}
\tightlist
\item
  Let A, B be two subrings of K and C the subring of K generated by A ∪ B. Show that the following conditions are equivalent:
\end{enumerate}

\((\alpha)\) For every local ring \(\mathbf{Q} \subset \mathbf{K}\) containing A and B, \(\mathbf{A}_{\mathfrak{p}} = \mathbf{B}_{\mathfrak{q}}\) , where \(\mathfrak{p} = \mathfrak{m}(\mathbf{Q}) \cap \mathbf{A}\) , \(\mathfrak{q} = \mathfrak{m}(\mathbf{Q}) \cap \mathbf{B}\) .

(β) For every prime ideal r of C, \(A_{p} = B_{q}\) , where \(p = r \cap A\) , \(q = r \cap B\) .

(γ) If M ∈ L(A) and N ∈ L(B) are related, they are identical.

(δ) L(A) ∩ L(B) = L(C).

(Use (a) and (b)).

\begin{enumerate}
\def\labelenumi{\arabic{enumi}.}
\setcounter{enumi}{1}
\item
  For an integral domain A to be a valuation ring, it is necessary and sufficient that every ideal of A be irreducible (Algebra, Chapter II, § 2, Exercise 16).
\item
  Show that every valuation ring is coherent (Chapter I, § 2, Exercise 12).
\item
  Let K be a field and F a set of valuation rings of K totally ordered by inclusion. Show that the intersection of the rings belonging to F is a valuation ring of K.
\item
  Let K be a field, A a integrally closed subdomain of K with K as a field of fractions and \((\mathbf{A}_{\alpha})\) a family of valuation rings of K whose intersection is A. Then if L is an extension of K, the integral closure of A in L is the intersection of the valuation rings of L which dominate one of the \(A_{\alpha}\) .
\end{enumerate}

¶ 6. Let R be a ring and A a subring of R such that S = R - A is non-empty and the product of two elements of S belongs to S.

\begin{enumerate}
\def\labelenumi{(\alph{enumi})}
\item
  Let \(a, a'\) be two elements of \(\mathbf{A}\) and \(s, s'\) two elements of \(S\) such that \(sa \in \mathbf{A}\) and \(s'a' \in \mathbf{A}\) . Show that one of the two elements \(sa', s'a\) belongs to \(\mathbf{A}\) . Deduce that the set \(p_1\) of elements \(a \in \mathbf{A}\) for which there exists \(s \in S\) such that \(sa \in \mathbf{A}\) is a prime ideal of \(\mathbf{A}\) .
\item
  Show that the set \(\mathfrak{p}_0\) of \(a \in \mathbf{A}\) such that \(sa \in \mathbf{A}\) for all \(s \in \mathbf{S}\) is a prime ideal of \(\mathbf{R}\) and \(\mathbf{A}\) ; it is the greatest ideal of \(\mathbf{R}\) contained in \(\mathbf{A}\) and \(\mathfrak{p}_0 \subset \mathfrak{p}_1\) .
\item
  The set \(\mathfrak{n}\) of elements \(x \in \mathbb{R}\) such that there exists \(s \in \mathbb{S}\) for which \(sx = 0\) is both an ideal of \(\mathbb{R}\) and an ideal of \(\mathbb{A}\) and \(\mathfrak{n} \subset \mathfrak{p}_0\) .
\item
  The ring \(\mathbf{A}\) is integrally closed in \(\mathbb{R}\) .
\item
  If \(\mathbf{R}\) is a field, \(\mathbf{A}\) is a valuation ring of \(\mathbf{R}\) , \(p_0 = (0)\) and \(p_1\) is the maximal ideal of \(\mathbf{A}\) .
\end{enumerate}

¶ 7. Let R be a ring. An ordered pair (A, p) consisting of a subring A of R and a prime ideal p of A is called maximal in R if there exists no ordered pair (A', p') consisting of a subring A' of R and a prime ideal p' of A' such that A ⊂ A', A' ≠ A and p' ∩ A = p.

\begin{enumerate}
\def\labelenumi{(\alph{enumi})}
\item
  For every subring B of R and every prime ideal q of B there exists an ordered pair \((A, p)\) , maximal in R, such that \(A \supset B\) and \(q = p \cap B\) .
\item
  For an ordered pair (A, p) to be maximal in R, it is necessary and sufficient that, for all \(s \in R - A\) , there exist a finite family of elements \(c_{i} \in p\) ( \(1 \leqslant i \leqslant n\) ) such that the element \(b = c_{1}s + c_{2}s^{2} + \cdots + c_{n}s^{n}\) belongs to A - p (use Chapter II, §2, no. 5, Corollary 3 to Proposition 11).
\item
  Show that, if the ordered pair \((\mathbf{A},\mathfrak{p})\) is maximal in \(\mathbf{R}\) , the product of two elements of \(\mathbf{R} - \mathbf{A}\) belongs to \(\mathbf{R} - \mathbf{A}\) . (If \(s,s'\) are in \(\mathbf{R} - \mathbf{A}\) and \(ss' = a\in \mathbf{A}\) , consider the least integers \(n,m\) such that there exists \(c_{i}\in \mathfrak{p}\) ( \(1\leqslant i\leqslant n\) ) for which \(b = \sum_{i}c_{i}s^{i}\in \mathbf{A} - \mathfrak{p}\) and \(c_{j}'\in \mathfrak{p}\) ( \(1\leqslant j\leqslant m\) ) for which
\end{enumerate}

\[
b ^ {\prime} = \sum_ {j} c _ {j} ^ {\prime} s ^ {\prime j} \in A - p;
\]

supposing for example \(n \geqslant m\) , consider the product \(bb' \in A - p\) and obtain a contradiction from the definition of the integer \(n\) .)

\begin{enumerate}
\def\labelenumi{(\alph{enumi})}
\setcounter{enumi}{3}
\item
  Under the hypotheses of (c), we write \(S = R - A\) ; show that the ideal \(p_1\) of \(A\) consisting of the \(a \in A\) such that there exists \(s \in S\) for which \(sa \in A\) is contained in \(p\) (use (b)). A fortiori the ideals \(p_0\) of \(A\) consisting of the \(a \in A\) such that \(sa \in A\) for all \(s \in S\) is contained in \(p\) (Exercise 6 (b)). We write \(A' = A / p_0\) , \(p' = p / p_0\) and \(R' = R / p_0\) ; the ordered pair \((A', p')\) is then maximal in the integral domain \(\mathbf{R}'\) and, for all non-zero \(a' \in \mathbf{A}'\) , there exists \(s' \in \mathbf{S}' = \mathbf{R}' - \mathbf{A}'\) such that \(s'a' \in \mathbf{S}'\) .
\item
  Let \((\mathbf{A},\mathfrak{p})\) be a maximal ordered pair in the integral domain \(\mathbf{R}\) such that, for all non-zero \(a\in \mathbf{A}\) , there exists \(s\in S = R - A\) such that \(sa\in S\) . We write \(\mathrm{S}_0 = \mathrm{S}\cup \{1\}\) and denote by \(\mathbb{R}_0\) the ring of fractions \(\mathrm{S}_0^{-1}\mathrm{R}\) ; the canonical homomorphism \(\mathbb{R}\to \mathbb{R}_0\) is injective and therefore identifies \(\mathbb{R}\) with a subring of \(\mathbb{R}_0\) . Show that \(\mathbb{R}_0\) is a field identified with the field of fractions of \(\mathbb{R}\) . Let \((\mathbf{B},\mathfrak{q})\) be a maximal ordered pair in \(\mathbb{R}_0\) such that \(\mathbf{B}\supset \mathbf{A}\) and \(\mathfrak{q}\cap \mathbf{A} = \mathfrak{p}\) . Show that \(\mathbf{B}\cap \mathbf{R} = \mathbf{A}\) ; in other words, \(\mathbf{A}\) is the intersection of \(\mathbf{R}\) and a valuation ring of \(\mathbb{R}_0\) and \(\mathfrak{p}\) is the intersection of \(\mathbf{A}\) and the maximal ideal of this valuation ring.
\item
\end{enumerate}

¶ 8. Given a ring A, a polynomial P(X \(_{1}\) , \ldots, X \(_{n}\) ) with coefficients in A is called dominated if it is of the form X \(^{\alpha}\) + ∑ \(_{\beta} c_{\beta}\) X \(^{\beta}\) , where β \textless{} α in the product ordered set N \(^{n}\) for all β such that c \(_{\beta}\) ≠ 0. In a ring R, a subring A is called paravaluative if, for every dominated polynomial P ∈ A{[}X \(_{1}\) , \ldots, X \(_{n}\) {]} (any n), P(s \(_{1}\) , \ldots, s \(_{n}\) ) ≠ 0 for all elements s \(_{i}\) ∈ R − A (1 ≤ i ≤ n).

\begin{enumerate}
\def\labelenumi{(\alph{enumi})}
\item
  Show that, if \(\mathbf{A}\) is paravaluative in \(\mathbb{R}\) , the product of two elements of \(\mathbf{S} = \mathbf{R} - \mathbf{A}\) belongs to \(\mathbf{S}\) .
\item
  Let A be a subring of a ring R such that the product of two elements of S = R - A belongs to S; let \(p_{0}\) be the prime ideal of A and R consisting of the \(a \in A\) such that \(sa \in A\) for all \(s \in S\) (Exercise 6 (b)). For A to be paravaluative in R, it is necessary and sufficient that \(A/p_{0}\) be paravaluative in \(R/p_{0}\) .
\item
  Let R be a ring and \(h: \mathbf{R} \to \mathbf{R}'\) a ring homomorphism. If \(\mathbf{A}'\) is a paravaluative subring of \(\mathbf{R}'\) , \(\mathbf{A} = \bar{h}^{-1}(\mathbf{A}')\) is paravaluative in R. Deduce that, if \((\mathbf{A}, p)\) is a maximal ordered pair in R (Exercise 7), A is paravaluative in R (use (b) and Exercise 7 (d) and (e)).
\item
  Let \(\mathbf{R}'\) be a ring and \(\mathbf{R}\) a subring of \(\mathbf{R}'\) . If \(\mathbf{A}\) is a paravaluative subring of \(\mathbf{R}\) , show that there exists a paravaluative subring \(\mathbf{A}'\) of \(\mathbf{R}'\) such that \(\mathbf{A}' \cap \mathbf{R} = \mathbf{A}\) . (If \(S = R - A\) and \(S_0 = S \cup \{1\}\) , consider the ring \(T' = S_0^{-1}R'\) and the canonical homomorphism \(h: R' \to T'\) . Let \(\mathbf{B}\) be the subring of \(T'\) generated by \(h(A) \cup (h(S))^{-1}\) and \(\mathfrak{b}\) the ideal of \(\mathbf{B}\) generated by \((h(S))^{-1}\) . Show that \(\mathfrak{b} \neq \mathbf{B}\) ; consider a maximal ordered pair \((\mathbf{A}^{\prime \prime}, \mathfrak{p}^{\prime \prime})\) in \(\mathbf{T}'\) such that \(\mathbf{B} \subset \mathbf{A}''\) and \(\mathfrak{b} \subset \mathfrak{p}''\) and take \(\mathbf{A}' = \overline{h}^{-1}(\mathbf{A}'')\) .
\end{enumerate}

Conclude from this that, if R is an integral domain, the paravaluative sub-rings of R are the intersections of R with the valuation rings of the field of fractions of R.

¶ 9. (a) Let R be a ring, A a subring of R, x an element of R and S the set of \(x^{k}\) ( \(k \geqslant 0\) ) in R. Show that, for x to be integral over A, it is necessary and sufficient that \(x/1 \in A[1/x]\) in the ring \(S^{-1}R\) . If x is not integral over A, show that there exists a maximal ideal m of \(A[1/x]\) containing 1/x and that the inverse image of m under the canonical homomorphism \(A \to A[1/x]\) is a maximal ideal of A.

\begin{enumerate}
\def\labelenumi{(\alph{enumi})}
\setcounter{enumi}{1}
\tightlist
\item
  Show that the integral closure A' of A in R is the intersection of the paravaluative subrings of R (Exercise 8) which contain A. (To see that A' is contained in each of these rings, use Exercise 8 (a) and Exercise 6 (d); to see that, if \(x \in R\) is not integral over A, there exists a subring B of R which is paravaluative in R and such that \(x \notin B\) , use (a), Exercise 7 (a) and Exercise 8 (c), arguing as in Theorem 3 of no. 3.)
\end{enumerate}

